\section{Dataset and Data Requirements}

All experiments use the ISIC-2019 training collection, a public set of 25{,}331
dermoscopic images labelled across the eight diagnostic classes listed in
Chapter~1. The collection is itself a union of three sources --- HAM10000,
BCN20000, and MSK --- which we infer from the image-ID prefix, because keeping
track of the source matters for building an honest train/test split. Each image
also carries clinical metadata: approximate age, sex, the general anatomical
site of the lesion, and a lesion identifier that groups multiple images of the
same physical lesion.

The defining property of this dataset, and the one that drives most of our
design decisions, is its class imbalance. Table~\ref{tab:classdist} shows the
per-class counts for the full collection. The largest class, nevus, is more than
fifty times larger than the smallest, dermatofibroma. Several of the rare
classes (for example SCC and AK) are clinically important, so a model that
simply learns the head of this distribution would be useless in practice. This
is the central reason we report balanced multi-class accuracy rather than overall
accuracy throughout.

\begin{table}[H]
\centering
\caption{Class distribution of the full ISIC-2019 collection (before
deduplication) and imbalance relative to the largest class (NV).}
\label{tab:classdist}
\begin{tabular}{llrr}
\toprule
Code & Full name & Images & Ratio vs.\ NV \\
\midrule
MEL  & Melanoma                  & 4{,}522  & 2.8$\times$  \\
NV   & Melanocytic nevus         & 12{,}875 & 1$\times$    \\
BCC  & Basal cell carcinoma      & 3{,}323  & 3.9$\times$  \\
AK   & Actinic keratosis         & 867      & 14.8$\times$ \\
BKL  & Benign keratosis          & 2{,}624  & 4.9$\times$  \\
DF   & Dermatofibroma            & 239      & 53.9$\times$ \\
VASC & Vascular lesion           & 253      & 50.9$\times$ \\
SCC  & Squamous cell carcinoma   & 628      & 20.5$\times$ \\
\bottomrule
\end{tabular}
\end{table}

The metadata is genuinely incomplete: roughly 30\% of age values, about 20\% of
sex values, and about 15\% of anatomical-site values are missing, with the
recorded ages centred near 54 years. Rather than fill these gaps with guessed
values, our requirement is that the model treat ``missing'' as a first-class,
learnable signal, which is reflected directly in the architecture (Chapter~4).

\section{System Requirements}

\subsection{Hardware}

The project targets a single-workstation setting rather than a data centre. All
training and evaluation run on one machine with an NVIDIA RTX 4070 Ti Super
(16~GB of VRAM), an Intel Core i7-14700 CPU, and 64~GB of system memory. The
16~GB memory ceiling is a real constraint: it sets the batch sizes we can use and
makes mixed-precision training a necessity rather than an option.

\subsection{Software}

The implementation is in PyTorch with the \texttt{timm} library supplying the
pretrained CNN backbones. The codebase is organised as importable modules
(preprocessing, datasets, losses, models, training, analysis) so that every
stage of the pipeline is reproducible from configuration files. Figures and the
results tables in this thesis are generated directly from the saved evaluation
outputs.

\section{Design Constraints}

Several hard constraints shaped the implementation. The models run under Windows,
where the PyTorch data loader is prone to deadlocks with multiple worker
processes; we therefore load data with zero worker processes and compensate by
decoding the entire dataset once into a memory-mapped array, after which each
batch is essentially a memory copy. The 16~GB GPU budget means automatic mixed
precision and channels-last memory layout are used everywhere, and the
larger models rely on gradient accumulation to reach a useful effective batch
size. Beyond the hardware, we imposed an explicit efficiency budget on the
lightweight model --- under 6M parameters and under 1 GMAC --- because a model
that cannot run on modest clinic hardware does not meet the goal of the project.

\section{Impacts}

The intended impact of this work is a decision-support tool that is small enough
to run where care is actually delivered, and that uses the same contextual
information a clinician would. A compact model lowers the barrier to deployment
on portable or low-cost hardware, which is most valuable exactly where
specialist coverage is thin. At the same time, the system is meant to assist and
not replace a clinician; it produces a ranked set of possibilities rather than a
final diagnosis. We are also conscious of the ethical dimension: the data is
de-identified and used only for research, the metadata is used solely as model
input, and the reliance on balanced accuracy is itself an ethical choice, since
it forces the model to take the rare, often malignant, classes seriously rather
than optimising a number that the common classes dominate.
